\chapter{Metodologia}

Per assolir els objectius d'aquest treball, s'ha dissenyat una metodologia mixta que combina la recerca documental amb el treball de camp i la creació d'un prototip físic. El procés s'ha dividit en quatre fases diferenciades:

\section{Fase 1: Recerca i fonamentació teòrica}
En aquesta etapa s'ha realitzat una revisió bibliogràfica extensa sobre la comunicació augmentativa i alternativa (CAA) i la teoria de la integració sensorial. L'objectiu ha estat entendre com els estímuls sensorials (tacte, llum, so) poden actuar com a pont per facilitar l'expressió de necessitats en persones amb dificultats severes.

\section{Fase 2: Treball de camp i observació}
S'ha dut a terme una observació directa en centres especialitzats, destacant la col·laboració amb el centre \textbf{Horitzó}. Mitjançant entrevistes a professionals de la logopèdia i l'educació especial, s'han detectat les mancances en els recursos actuals i les oportunitats que ofereix una eina tàctil i interactiva.

\section{Fase 3: Disseny i construcció de la taula sensorial}
Aquesta és la part pràctica central del treball. La metodologia de construcció ha seguit aquests passos:

\subsection{Materials i components}
S'ha realitzat una selecció acurada dels components per garantir una experiència sensorial completa i segura. La següent taula resumeix els elements principals utilitzats:

\begin{center}
    \renewcommand{\arraystretch}{1.5}
    \rowcolors{2}{SoftBlue!10}{white}
    \begin{tabularx}{\textwidth}{l X l}
        \rowcolor{SoftBlue!30}
        \textbf{Component} & \textbf{Descripció i Utilitat} & \textbf{Tipus d'estímul} \\ \toprule
        Polsadors LED & Botons grans que s'il·luminen al tacte. & Visual / Tàctil \\
        Panells de textura & Superfícies de silicona, fusta i vellut. & Tàctil \\
        Mòdul de veu & Reproducció de missatges pre-gravats. & Auditiu \\
        Fibra òptica & Filaments de llum per captar l'atenció. & Visual \\
        Carcassa de fusta & Estructura base ergonòmica i resistent. & Estructural \\ \bottomrule
    \end{tabularx}
\end{center}

\subsection{Procés de muntatge}
La metodologia de construcció ha seguit aquests passos:

\section{Fase 4: Validació i recollida de dades}
Un cop construïda la taula, s'han realitzat proves pilot amb usuaris per observar-ne la reacció. S'ha utilitzat una graella d'observació per recollir dades sobre:
\begin{itemize}
    \item Temps de resposta a l'estímul.
    \item Nivell d'interès i participació.
    \item Capacitat de l'usuari per associar l'estímul amb una acció comunicativa.
\end{itemize}
